\section{Additional Method Details}
\label{sec:additional-method}

This appendix makes explicit the contracts that connect the three training
stages.  This is important because the central claim of AMPT is not that a
trajectory with a high offline score is automatically a useful target, but
that supervision, rollout credit, and policy-relative refinement must agree on
what the current policy can safely learn.  We use three provenance labels
throughout.  \emph{Paper-reported} denotes a value stated in the main paper,
\emph{scene-level reproduced} denotes a value recomputed from an archived
per-scene evaluator export, and \emph{proposed reproduction} denotes a fully
specified setting for rerunning an experiment whose resolved training manifest
was not retained.  The last category is a protocol, not an observed result.

\subsection{Policy-Compatible Multi-Trajectory Supervision}
\label{sec:pc-mts-details}

\paragraph{Candidate construction and source balancing.}
For a scene $s$, let $\tau_s^0$ be the logged trajectory and let
$\mathcal C_s=\cup_m\mathcal C_s^{(m)}$ be the union of candidates supplied by
the configured proposal generators.  PC-MTS is source agnostic: a source may
be another policy seed, a structured trajectory expander, or a specialist
planner, but source identity is retained in the manifest and geometric near-
duplicates are merged before scoring.  The logged target is immutable and
is never removed.  A final Stage-1 source-count manifest was not found.  To
make the controlled reconstruction executable without presenting invented
counts as measurements, our \emph{proposed reproduction} caps every non-GT
source at eight candidates per scene before filtering.  When comparing Score,
Pareto, and PC-MTS, we use the same source
cap and retain the same number of non-GT candidates by taking the best valid
items under the corresponding rule.  Thus a difference cannot be attributed
to exposing one variant to a larger target set.

\begin{table*}[t]
\centering
\small
\caption{Candidate-source and audit budgets for PC-MTS on the NAVSIM training universe. Counts tagged proposed are reproduction settings rather than observations; higher/lower is not applicable.}
\label{tab:candidate_sources}
\resizebox{\textwidth}{!}{%
\begin{tabular}{lccccc}
\toprule
source & role & per\_scene\_budget & scene\_universe & retention\_or\_use & status \\
\midrule
Logged trajectory & immutable fallback and retention target & 1 & 103,288 expected navtrain scenes & always retained & paper-described / historical contract \\
Frozen-policy rollout bank & compatibility support estimate & 16 & 103,288-scene reproduction universe & distance calibration only & proposed reproduction setting \\
Multi-source candidate proposals & quality and Pareto targets & 8 per source & 103,288-scene reproduction universe & subject to quality/Pareto/compatibility gates & proposed reproduction budget \\
Correlated local perturbations & local feasibility audit & 8 per retained proposal & same candidate scenes & pass-rate threshold 0.75 & proposed reproduction setting \\
\bottomrule
\end{tabular}%
}
\end{table*}


\paragraph{Quality, feasibility, and Pareto filtering.}
Every complete trajectory is scored by the same offline evaluator used for
training diagnostics.  We first require hard feasibility, then a sufficient
aggregate score, and only then form the scene-wise component Pareto front.  A
candidate $a$ dominates $b$ only if it is no worse on every active Pareto
objective and is strictly better on at least one.  This order matters: a high
progress value cannot compensate for collision or leaving the drivable area.
The paper does not archive the numerical Stage-1 aggregate cutoff.  In the
\emph{proposed reproduction}, a training-split quantile is selected once so
that comparison variants are matched to the minimum retained count, followed
by a source-balanced top-$k$ tie break.  The sealed cutoff is applied
identically to all splits and is never used to synthesize a reported score.

The benchmark-specific partition is summarized below.  On NAVSIM v1, NC and
DAC are hard gates; DDC is protected; and EP, TTC, and comfort are optimized
inside the admissible region.  On NAVSIM v2, NC, DAC, DDC, and TLC are the
multiplicative compliance gates.  The historical metric adapter uses EP, TTC,
and a composite $Q=(\mathrm{LK}+\mathrm{HC}+\mathrm{EC})/3$ as its three
continuous Pareto axes; the official EPDMS report still exposes LK, HC, and EC
separately.  For binary scorer outputs the \emph{proposed reproduction}
requires a value of one up to $10^{-6}$.  The archived safety contract uses
reference-relative tolerances of $0.01$ for DDC/comfort and $0.02$ for TTC/EP;
TLC remains a hard v2 gate.  These are historical implementation values, not a
claim about the unresolved configuration of the paper checkpoint.

\begin{table}[t]
\centering
\small
\caption{NAVSIM v1/v2 metric partition used by PC-MTS and FF-PGRPO. Hard feasibility is tested before protected non-regression and Pareto quality.}
\label{tab:metric_partition}
\resizebox{\linewidth}{!}{%
\begin{tabular}{lcccc}
\toprule
benchmark & hard\_feasibility & protected\_guards & pareto\_quality & aggregate \\
\midrule
NAVSIM v1 & NC, DAC & DDC; non-regression in TTC and feasibility & EP, TTC, C & PDMS \\
NAVSIM v2 & NC, DAC, DDC, TLC & DDC/TLC feasibility; non-regression in TTC & EP, TTC, LK, HC, EC & EPDMS \\
\bottomrule
\end{tabular}%
}
\end{table}


\paragraph{Policy compatibility.}
Let $\mathcal R_s=\{\rho_j\}_{j=1}^{M}$ be rollouts sampled from the frozen
GT-only policy under the same command as the candidate.  We standardize
waypoint displacement with per-timestep, command-conditioned robust scales
fitted on the calibration split, and write the resulting trajectory distance
as $d(\tau,\rho)$.  The archived implementation additionally exposes final
displacement, heading, and velocity diagnostics, but the proposed controlled
comparison changes only the normalized waypoint statistic.  The compatibility
statistic is
\begin{equation}
 d_{\mathrm{PC}}(\tau,\mathcal R_s)
 =\frac{1}{K}\sum_{\rho\in\operatorname{NN}_{K}(\tau;\mathcal R_s)}
 d(\tau,\rho).
\end{equation}
The threshold $q_c$ is the command-specific quantile of leave-one-out policy
rollout distances on calibration scenes.  The archived code proves
leave-one-out calibration within each policy cloud, but its retained manifest
does not prove driving-log-level independence.  The proposed reconstruction
therefore enforces log-disjoint calibration and query scenes.  The paper states the general
$K$-nearest statistic but does not report $M$, $K$, or $q_c$.  The archived
implementation realizes the valid $K=1$ special case.  Our \emph{proposed
reproduction} uses $M=16$, $K=3$, and $q_c$ at the 95th percentile, with
command groups merged into a global fallback only when a command group has
fewer than 16 calibration observations.  This choice is also the center point
of the supplied sensitivity grid, rather than a retrospectively selected
result.

\paragraph{Local feasibility tube.}
Compatibility to one rollout neighborhood does not by itself establish that a
target is robust to small prediction errors.  For every candidate that passes
the previous gates, PC-MTS adds correlated perturbations to $(x,y,\psi)$ and
the implied velocity sequence.  The perturbations are sampled from normalized
policy residuals, smoothed over time, and projected back to the command-consistent
trajectory representation before exact re-evaluation.  This changes whole
trajectory segments rather than independently jittering waypoints.  The
archived implementation supports 16 or 32 perturbations; the \emph{proposed
reproduction} uses eight lower-cost perturbations, preserves the empirical
temporal covariance of policy residuals, and accepts a candidate when at least
six of eight perturbations pass
the hard feasibility gates and the protected guards.  The P1 configuration
evaluates 4/8/16 perturbations and does not tune the threshold on the test
split.

\paragraph{Target sampling and fallback.}
After filtering, the training set for scene $s$ is
$\mathcal T_s=\{\tau_s^0\}\cup\mathcal A_s$, where $\mathcal A_s$ contains
accepted candidates.  A training update draws exactly one target for the
scene, so scenes with many proposals do not receive larger gradients.  The
\emph{proposed reproduction} draws the logged target with probability $0.5$
and otherwise samples uniformly from $\mathcal A_s$; if $\mathcal A_s$ is
empty, it deterministically uses $\tau_s^0$.  Candidate identity is sampled
before the diffusion timestep, so all noisy versions in that update correspond
to a single coherent target.  Fine-tuning uses the same denoising objective as
GT-only training and produces $\pi_{\mathrm{initial}}$.

\paragraph{Implementation correspondence.}
The source-balanced builder and immutable-GT path are implemented at
\path{41a8613:navsim/agents/recogdrive/curriculum/builder.py}
and \path{41a8613:navsim/agents/recogdrive/curriculum/proposals.py}.
Physical distance and calibrated reachability are at
\path{41a8613:navsim/agents/recogdrive/curriculum/reachability.py};
the feasibility and correlated-perturbation checks are at
\path{41a8613:navsim/agents/recogdrive/curriculum/safety.py} and
\path{41a8613:navsim/agents/recogdrive/curriculum/robustness.py}.
These are historical mechanism-level implementations; the audit does not claim
that this commit is a complete hash-identical release of the final checkpoint.

\subsection{Feasibility-First Pareto GRPO}
\label{sec:ff-pgrpo-details}

\paragraph{A coherent scene reference.}
FF-PGRPO compares a rollout with one executable reference trajectory, never
with an envelope assembled from different trajectories.  The ordered source
set is: the logged trajectory, a retained PC-MTS candidate, and the current
single-trajectory prediction.  The paper-level rule selects the highest-scoring
feasible row.  The audited historical selector makes this precise by ranking
validity, hard feasibility, reference guard, robustness, core eligibility,
PDMS, distance to the current prediction, and source priority lexicographically.
Thus no higher score can bypass an earlier safety/eligibility tier.  If none is feasible, the
logged trajectory is retained only as a weak recovery target when it passes
the evaluator's basic physical checks.  Otherwise no recovery target is used
and every sampled rollout receives non-positive credit.  A component-wise
envelope is invalid because its NC, progress, and TTC values may come from
mutually incompatible plans; positive credit relative to such a synthetic
point would have no executable safety witness.

\paragraph{Group states and admissibility.}
For each scene the paper reports a group of $G=16$ policy rollouts.  We assign
the group to one of three states.  An \emph{all-feasible} group contains only
rollouts that satisfy hard feasibility; its quality statistic normalizes the
continuous score components because the binary gates are constant.  A
\emph{mixed} group contains both feasible and infeasible rollouts; feasibility
is applied first and feasible trajectories must additionally pass the coherent
reference guards.  An \emph{all-infeasible} group contains no feasible
rollout; it is not allowed to create a positive advantage by merely ranking
unsafe samples against one another.

For NAVSIM v1, admissibility requires NC and DAC and protects DDC, TTC, and an
EP floor relative to the coherent reference.  NAVSIM v2 additionally treats
DDC and TLC as hard compliance checks.  The archived safety-contract defaults
use $\mathrm{DDC}(\tau)\geq\mathrm{DDC}(r)-0.01$,
$\mathrm{TTC}(\tau)\geq\mathrm{TTC}(r)-0.02$, and
$\mathrm{EP}(\tau)\geq\mathrm{EP}(r)-0.02$.  The joint EP--TTC guard prevents
the optimizer from receiving positive credit for increasing safety merely by
stopping, or for increasing progress by collapsing the collision margin.

\paragraph{Credit assignment.}
Within a group that contains a feasible rollout, let $\widehat A_i$ be a
quality-shaped, group-relative score and let $P_i$ indicate membership in the
admissible Pareto front.  FF-PGRPO assigns
\begin{equation}
 A_i=\begin{cases}
 \max(\widehat A_i,0),
   & \substack{\text{feasible, guard-passing,}\\\text{and Pareto-positive}},\\
 \min(\widehat A_i,0),
   & \substack{\text{feasible but not}\\\text{admissible Pareto}},\\
 -\lambda_{\mathrm{unsafe}}+\min(\widehat A_i,0),
   & \text{infeasible}.
 \end{cases}
\end{equation}
Thus aggregate quality controls magnitude, whereas feasibility, a coherent
reference, and non-domination control the sign.  For mixed groups the quality
term also includes the EP-floor and joint EP--TTC penalties.  The historical
implementation preserves signs with RMS scaling rather than mean centering and
uses an archived default clip of $[-3,3]$; the \emph{proposed reproduction}
uses the predeclared sensitivity setting $[-5,5]$.

For an all-infeasible group, violation severity is the normalized sum of hard
gate deficits and protected-metric deficits.  Credit is the non-positive
within-group rank of negative severity, shifted by the unsafe penalty.  The
group loss is downweighted to $0.25$; when a valid recovery reference exists,
a diffusion recovery term with coefficient $0.25$ is added.  The proposed
regularization uses KL weight $0.02$ and anneals the BC coefficient linearly
from $0.10$ to $0.05$ over the first five epochs.  These numerical values are
historical implementation defaults and proposed rerun values, not recovered
resolved parameters for the paper's final FF-PGRPO run.

\paragraph{Diffusion credit and inference.}
One trajectory-level advantage is broadcast to every transition in that
trajectory's denoising chain.  The audited planner uses a discounted
transition-wise mean, while the paper objective writes the chain likelihood in
accumulated form; the proposed reconstruction preserves the planner reduction
and logs it in the resolved manifest.  This avoids assigning mutually
inconsistent metric credit to intermediate noisy states that are not executable
plans.  At inference, AMPT runs one denoising chain and emits one trajectory.
Group sampling, the Pareto computation, the reference selector, the offline
evaluator, and candidate reranking are absent from test time.

\paragraph{Implementation correspondence.}
Group classification, Pareto gating, violation ranking, sign-preserving
normalization, and clipping are at
\path{41a8613:navsim/agents/recogdrive/stage3_ff_pareto.py}.
The coherent-row API is at
\path{41a8613:navsim/agents/recogdrive/stage3_reference_selector.py},
and positive-credit assertions are at
\path{41a8613:navsim/agents/recogdrive/stage3_positive_credit_contract.py}.
The archived tensor layout is anchor-local; the paper-level protocol is the
union of 16 scene rollouts.  Because a final resolved layout manifest was not
found, we do not infer the number of anchors from a class default.

\subsection{Adaptive Pareto-Guided Policy Refinement}
\label{sec:apr-details}

\paragraph{Dynamic multi-source pool.}
At APR round $r$, the current policy $\pi_r$ supplies one local prediction for
every scene.  Candidate teachers come from historical policy checkpoints,
independent GRPO seeds, structured trajectory expansion, and safety-,
progress-, and regularity-oriented specialists.  Source labels survive
deduplication so per-source acceptance can be audited, but final three-round
source proportions were not archived.  The audited initial construction keeps
its observed 25,402/13,859 source mixture.  The proposed source-ablation reports
each source separately instead of claiming an unrecorded balancing rule.

\paragraph{Teacher audit.}
For a policy prediction $p_s$ and candidate teacher $t$, APR first requires
feasibility and an aggregate margin $R(t)-R(p_s)\geq m_R$.  It then applies
absolute and reference-relative protected guards.  Behavioral distance,
curvature, and jerk rank otherwise admissible teachers; they are preferences,
not claimed historical hard thresholds.  In the audited APR branch,
$m_R=0$ with a strict positive comparison, DDC has absolute floor 0.99 and
maximum drop 0.01, TTC has absolute floor 0.95 and maximum drop 0.02, and the
jerk coefficient 0.005 is a soft training regularizer.  At the interpolation
audit, the proposed reconstruction additionally applies the command-conditioned
95th-percentile compatibility calibration from PC-MTS.

\paragraph{Local interpolation and exact re-evaluation.}
An accepted teacher is not copied directly.  APR evaluates
\begin{equation}
 \tau_{s,\alpha}=p_s+\alpha(t-p_s)
\end{equation}
for a descending list of interpolation coefficients.  The first (largest)
coefficient that remains feasible, improves the aggregate score, preserves
protected metrics, and passes the current compatibility test becomes the local
target.  Every interpolated trajectory is scored anew; endpoint scores are
never linearly interpolated.  One audited branch used $\alpha=0.25$.  The
\emph{proposed reproduction} fixes the descending list to
$[1.0,0.75,0.50,0.25]$, which includes that audited setting and tests the
teacher endpoint before progressively more conservative steps.  If all
coefficients fail, the teacher is
rejected for the current round rather than kept with a smaller unscored step.

\paragraph{Weighted distillation, retention, and lifecycle.}
For verified gain $\Delta_s=R(\tau_{s,\alpha})-R(p_s)>0$, the target weight is
\begin{equation}
 w_s=\min\{5,\exp(\Delta_s/0.05)\}.
\end{equation}
The temperature $0.05$ and clip 5 are present in an audited APR training
branch.  The same diffusion denoising loss is applied to weighted local
targets, while current predictions provide retention targets with audited
weight 1.0.  At most the top one positive-advantage target is used per scene; a
scene without a valid teacher contributes retention only.  After training, the
candidate checkpoint is promoted only if aggregate validation score increases
and none of the protected aggregate components regresses beyond the audit
tolerances.  The promoted checkpoint becomes the dynamic anchor, its
predictions and rollout bank are regenerated, and the complete teacher pool is
audited again.  A previously accepted teacher is defined as \emph{expired} when
it loses its margin or guard, while a previously rejected one is \emph{newly
activated} when it passes under the new policy.  Historical scripts rebuild
round files but do not persist these labels; the proposed pipeline obtains them
by joining teacher hashes across consecutive round manifests.  The paper reports three
rounds.  The proposed stop rule is the earlier of three promotions, an empty
teacher set, or no checkpoint passing the promotion rule; it does not turn a
small unverified score fluctuation into a promotion.

\paragraph{Optional branch consolidation.}
Safety, progress, and regularity branches may be consolidated after their
teacher updates.  This is an optional checkpoint operation, not a substitute
for teacher auditing.  The audited anchor-constrained variant keeps the top
20\% salient coordinates of each base-relative task vector, removes
coordinates with inconsistent directions, clips every branch displacement to
the norm of its accepted branch, and caps total displacement at $1.05$ times
the largest accepted branch displacement.  The anchor is the current promoted
policy and is therefore updated across rounds.  If no merge candidate passes
the same evaluation guard as an ordinary APR checkpoint, the best accepted
single branch is retained.

\paragraph{Implementation correspondence.}
Teacher audits are implemented at
\path{fb96a03:scripts/stage3/build_pdms_strict_teacher_buffer.py}
and \path{fb96a03:scripts/stage3/build_pdms_safety_repair_teacher_buffer.py}.
Interpolation and exact rescoring are at
\path{ab9124c:scripts/stage3/build_pdms_blended_teacher_reference.py},
with local safety-tube checks at
\path{ab9124c:scripts/stage3/filter_pdms_teacher_reference_by_safety_tube.py}.
Retention schedules are built by
\path{fb96a03:scripts/stage3/build_pdms_teacher_refinement_schedule.py},
and the optional merge is at
\path{64f5531:scripts/stage3/merge_conflict_free_anchored_task_vector.py}.
APR is an orchestrated sequence of these audited operations; the repository
does not contain a single monolithic \texttt{APR.run} entry point.
